La electrificación del transporte público contribuye a reducir las emisiones y el ruido en las ciudades, y Santiago ha avanzado hasta operar una de las flotas de buses eléctricos más grandes fuera de China \cite{worldbank2020chile}. Sin embargo, estos buses tienen autonomía limitada y requieren tiempo de recarga en electroterminales con capacidad restringida. Por ello, la asignación de buses a viajes debe coordinarse con las decisiones de cuándo y dónde cargar.

Una mala planificación puede aumentar la flota necesaria, el gasto energético y los kilómetros sin pasajeros, o generar planes incompatibles con la infraestructura disponible. Esto ocurre porque la demanda de viajes coincide parcialmente con tarifas eléctricas elevadas y las recargas se concentran al regresar a los electroterminales. Aprovechar cada bus para atender distintas rutas (\textit{interlining}) y distribuir las recargas durante el día ofrece oportunidades para reducir costos y presión sobre la infraestructura, contribuyendo a una operación económicamente sostenible.

Este informe continúa el Informe 1, que caracterizó el problema de asignación de vehículos eléctricos, conocido como E-VSP (\textit{Electric Vehicle Scheduling Problem}), y propuso una descomposición espaciotemporal con horizonte rodante. Dada la complejidad de implementación señalada en la revisión, esta entrega replantea la metodología y profundiza en las estrategias de agrupamiento de rutas.

El objetivo general es proponer y evaluar una metodología de planificación diaria de buses eléctricos que minimice el costo total de operación frente a un caso base. Para ello, se busca formular matemáticamente el problema, cuantificar el ahorro de compartir buses entre rutas, evaluar el efecto de la batería y los puestos de carga sobre la flota y el costo, estimar el beneficio de programar las recargas frente a una política reactiva e identificar los supuestos con mayor impacto en los resultados.

El entregable esperado es un plan operacional diario de asignación de buses, traslados sin pasajeros y recargas. Incluye las jornadas de cada bus, entendidas como la secuencia de expediciones y traslados entre su salida y regreso al electroterminal, y un programa que especifica cuándo y dónde cargar y cuánta energía reponer. Se complementa con el desglose del costo total y los indicadores de desempeño frente al caso base.
